\subsection{Bidirectional Data Flow Pipelines: Generator-Based Coroutines}

A generator normally sends values outward through \texttt{yield}. But a suspended generator can also receive values from the outside through \texttt{.send()}. This turns the generator into a simple coroutine: execution can pause, the caller can inject data, and the generator can continue using that injected data.

This is not preemptive multitasking. The generator is not interrupted by the operating system at an arbitrary instruction. It voluntarily suspends at \texttt{yield}, and the caller resumes it explicitly.

The basic direction changes from this:

\begin{verbatim}
generator -> caller
\end{verbatim}

to this:

\begin{verbatim}
generator -> caller
generator <- caller
\end{verbatim}

The same suspended frame is reused across both directions.

\subsubsection{Consumers and Transformers: Feeding In-Flight Data via the \texttt{.send()} Interface}

A \texttt{yield} expression can produce a value outward and later receive a value inward.

\begin{verbatim}
def receiver():
    print("receiver started")

    value = yield "ready"

    print(f"received: {value!r}")

gen = receiver()

first = next(gen)
print(f"first yielded value: {first!r}")

try:
    gen.send("D'oh!")
except StopIteration:
    print("receiver finished")
\end{verbatim}

Possible output:

\begin{verbatim}
receiver started
first yielded value: 'ready'
received: "D'oh!"
receiver finished
\end{verbatim}

The line:

\begin{verbatim}
value = yield "ready"
\end{verbatim}

has two sides:

\begin{itemize}
    \item \texttt{yield "ready"} sends \texttt{"ready"} outward and suspends;
    \item when the generator resumes, the value sent by \texttt{.send()} becomes the result of the \texttt{yield} expression and is assigned to \texttt{value}.
\end{itemize}

A generator must be started before it can receive a non-\texttt{None} value.

\begin{verbatim}
def receiver():
    value = yield "ready"
    print(f"received: {value!r}")

gen = receiver()

try:
    gen.send("D'oh!")
except TypeError as err:
    print(f"error type: {type(err)}")
    print(f"message:    {err}")
\end{verbatim}

Possible output:

\begin{verbatim}
error type: <class 'TypeError'>
message:    can't send non-None value to a just-started generator
\end{verbatim}

A just-created generator is positioned before the first line of its body. There is no suspended \texttt{yield} expression ready to receive the value yet.

The usual start operation is called \textit{priming} the generator.

\begin{verbatim}
def receiver():
    print("starting")
    value = yield "ready"
    print(f"received: {value!r}")

gen = receiver()

print("priming")
msg = next(gen)
print(f"priming result: {msg!r}")

print("sending")
try:
    gen.send("No soup for you!")
except StopIteration:
    print("done")
\end{verbatim}

Possible output:

\begin{verbatim}
priming
starting
priming result: 'ready'
sending
received: 'No soup for you!'
done
\end{verbatim}

Calling \texttt{next(gen)} is equivalent to \texttt{gen.send(None)} for this initial advance.

\begin{verbatim}
def receiver():
    print("starting")
    value = yield "ready"
    print(f"received: {value!r}")

gen = receiver()

msg = gen.send(None)
print(f"msg: {msg!r}")

try:
    gen.send("Ni!")
except StopIteration:
    print("done")
\end{verbatim}

Possible output:

\begin{verbatim}
starting
msg: 'ready'
received: 'Ni!'
done
\end{verbatim}

The generator object exposes the \texttt{send} method.

\begin{verbatim}
def receiver():
    value = yield "ready"
    print(value)

gen = receiver()

print(f"type(gen): {type(gen)}")
print()

selected_names = [
    "__next__",
    "send",
    "throw",
    "close",
    "gi_frame",
    "gi_code",
    "gi_running",
]

for name in selected_names:
    print(f"{name:<14} {name in dir(gen)}")
\end{verbatim}

Possible output:

\begin{verbatim}
type(gen): <class 'generator'>

__next__       True
send           True
throw          True
close          True
gi_frame       True
gi_code        True
gi_running     True
\end{verbatim}

A generator-based consumer can repeatedly receive values.

\begin{verbatim}
def printer():
    print("printer ready")

    while True:
        text = yield
        print(f"got: {text!r}")

gen = printer()

next(gen)

gen.send("D'oh!")
gen.send("No soup for you!")
gen.send("The answer is 42.")
\end{verbatim}

Possible output:

\begin{verbatim}
printer ready
got: "D'oh!"
got: 'No soup for you!'
got: 'The answer is 42.'
\end{verbatim}

The bare \texttt{yield} sends \texttt{None} outward and suspends. When resumed by \texttt{send(value)}, the \texttt{yield} expression evaluates to the sent value.

This consumer never naturally finishes because the loop is infinite. It waits at \texttt{yield}, receives a value, processes it, then loops back and waits again.

A transformer can receive a value and yield a transformed result.

\begin{verbatim}
def upper_transformer():
    print("transformer ready")

    while True:
        text = yield "waiting"
        result = text.upper()
        yield result

gen = upper_transformer()

print(next(gen))
print(gen.send("d'oh!"))
print(next(gen))
print(gen.send("no soup for you!"))
\end{verbatim}

Possible output:

\begin{verbatim}
transformer ready
waiting
D'OH!
waiting
NO SOUP FOR YOU!
\end{verbatim}

This generator has two suspension points inside the loop:

\begin{verbatim}
text = yield "waiting"
    sends "waiting"
    waits for input

yield result
    sends transformed output
    waits to be resumed again
\end{verbatim}

The caller must follow that rhythm. After receiving a result, the caller advances once more to return the generator to the input-waiting point.

A slightly clearer design yields only after receiving input.

\begin{verbatim}
def upper_transformer():
    text = yield "ready"

    while True:
        result = text.upper()
        text = yield result

gen = upper_transformer()

print(next(gen))
print(gen.send("d'oh!"))
print(gen.send("ni!"))
print(gen.send("the answer is 42."))
\end{verbatim}

Possible output:

\begin{verbatim}
ready
D'OH!
NI!
THE ANSWER IS 42.
\end{verbatim}

Here, each \texttt{send()} both resumes the generator with input and receives the transformed output.

The statement:

\begin{verbatim}
text = yield result
\end{verbatim}

means:

\begin{verbatim}
send result outward
pause
when resumed, store incoming value in text
\end{verbatim}

The local state survives between sends.

\begin{verbatim}
def counted_receiver():
    count = 0

    while True:
        text = yield f"ready count={count}"
        count = count + 1
        print(f"{count:>2}. {text}")

gen = counted_receiver()

print(next(gen))
print(gen.send("D'oh!"))
print(gen.send("No soup for you!"))
print(gen.send("Ni!"))
\end{verbatim}

Possible output:

\begin{verbatim}
ready count=0
 1. D'oh!
ready count=1
 2. No soup for you!
ready count=2
 3. Ni!
ready count=3
\end{verbatim}

The local variable \texttt{count} is preserved in the suspended generator frame.

The frame can be inspected.

\begin{verbatim}
def counted_receiver():
    count = 0

    while True:
        text = yield f"ready count={count}"
        count = count + 1
        print(f"{count:>2}. {text}")

gen = counted_receiver()

print(next(gen))
print(f"locals after prime: {gen.gi_frame.f_locals}")

print(gen.send("D'oh!"))
print(f"locals after send:  {gen.gi_frame.f_locals}")
\end{verbatim}

Possible output:

\begin{verbatim}
ready count=0
locals after prime: {'count': 0}
 1. D'oh!
ready count=1
locals after send:  {'count': 1, 'text': "D'oh!"}
\end{verbatim}

The generator frame stores the live local names. The generator is not restarted for every sent value.

A pipeline can connect a producer, transformer, and consumer manually.

\begin{verbatim}
def clean_transformer():
    text = yield "ready"

    while True:
        cleaned = text.strip().upper()
        text = yield cleaned

cleaner = clean_transformer()

next(cleaner)

items = ["  d'oh!  ", "  no soup for you!  ", "  ni!  "]

for item in items:
    result = cleaner.send(item)
    print(f"{item!r} -> {result!r}")
\end{verbatim}

Possible output:

\begin{verbatim}
"  d'oh!  " -> "D'OH!"
'  no soup for you!  ' -> 'NO SOUP FOR YOU!'
'  ni!  ' -> 'NI!'
\end{verbatim}

This is a bidirectional flow:

\begin{verbatim}
caller sends raw text into generator
generator sends cleaned text back to caller
\end{verbatim}

The generator is a suspended transformation machine.

\subsubsection{Exception Injection with \texttt{.throw()} and Controlled Shutdown with \texttt{.close()}}

The \texttt{.throw()} method injects an exception into the generator at its current suspension point.

\begin{verbatim}
def receiver():
    try:
        while True:
            text = yield "ready"
            print(f"got: {text}")
    except ValueError as err:
        print(f"receiver caught: {err}")
        yield "recovered"

gen = receiver()

print(next(gen))
print(gen.send("D'oh!"))
print(gen.throw(ValueError("bad input")))
\end{verbatim}

Possible output:

\begin{verbatim}
ready
got: D'oh!
ready
receiver caught: bad input
recovered
\end{verbatim}

The exception is raised inside the generator at the suspended \texttt{yield}. The generator may catch it with \texttt{try}/\texttt{except}.

The exception object is ordinary.

\begin{verbatim}
err = ValueError("bad input")

print(f"err:       {err}")
print(f"type(err): {type(err)}")
print()

selected_names = [
    "args",
    "__class__",
    "__str__",
    "__repr__",
    "with_traceback",
]

for name in selected_names:
    print(f"{name:<18} {name in dir(err)}")

print()
print(f"err.args: {err.args}")
\end{verbatim}

Possible output:

\begin{verbatim}
err:       bad input
type(err): <class 'ValueError'>

args               True
__class__          True
__str__            True
__repr__           True
with_traceback     True

err.args: ('bad input',)
\end{verbatim}

If the generator does not catch the injected exception, it propagates back to the caller.

\begin{verbatim}
def receiver():
    while True:
        text = yield "ready"
        print(f"got: {text}")

gen = receiver()

print(next(gen))

try:
    gen.throw(ValueError("bad input"))
except ValueError as err:
    print(f"caller caught: {err}")
\end{verbatim}

Possible output:

\begin{verbatim}
ready
caller caught: bad input
\end{verbatim}

The generator is no longer normally resumable after an uncaught exception terminates it.

\begin{verbatim}
def receiver():
    while True:
        text = yield "ready"
        print(f"got: {text}")

gen = receiver()

print(next(gen))

try:
    gen.throw(ValueError("bad input"))
except ValueError:
    print("uncaught exception terminated generator")

try:
    next(gen)
except StopIteration:
    print("generator is exhausted")
\end{verbatim}

Possible output:

\begin{verbatim}
ready
uncaught exception terminated generator
generator is exhausted
\end{verbatim}

The \texttt{.close()} method requests generator shutdown. Internally, it raises \texttt{GeneratorExit} at the suspension point.

\begin{verbatim}
def receiver():
    try:
        while True:
            text = yield "ready"
            print(f"got: {text}")
    finally:
        print("receiver cleanup")

gen = receiver()

print(next(gen))
print(gen.send("D'oh!"))

gen.close()

print("closed")
\end{verbatim}

Possible output:

\begin{verbatim}
ready
got: D'oh!
ready
receiver cleanup
closed
\end{verbatim}

The \texttt{finally} block runs during shutdown. This is the right place to release resources or print cleanup messages in educational examples.

After closing, the generator is exhausted.

\begin{verbatim}
def receiver():
    try:
        while True:
            text = yield "ready"
            print(f"got: {text}")
    finally:
        print("cleanup")

gen = receiver()

print(next(gen))
gen.close()

try:
    next(gen)
except StopIteration:
    print("cannot resume closed generator")
\end{verbatim}

Possible output:

\begin{verbatim}
ready
cleanup
cannot resume closed generator
\end{verbatim}

The generator object still exists as a Python object, but its suspended execution frame is gone.

\begin{verbatim}
def receiver():
    try:
        yield "ready"
    finally:
        print("cleanup")

gen = receiver()

print(f"before prime frame: {gen.gi_frame}")

print(next(gen))

print(f"before close frame: {gen.gi_frame}")

gen.close()

print(f"after close frame:  {gen.gi_frame}")
\end{verbatim}

Possible output:

\begin{verbatim}
before prime frame: <frame at 0x..., file "...", line 1, code receiver>
ready
before close frame: <frame at 0x..., file "...", line 3, code receiver>
cleanup
after close frame:  None
\end{verbatim}

Closing destroys the suspended frame state.

A generator should not yield a normal value while handling \texttt{GeneratorExit}. That violates the shutdown contract.

\begin{verbatim}
def bad_receiver():
    try:
        yield "ready"
    except GeneratorExit:
        yield "bad idea"

gen = bad_receiver()

print(next(gen))

try:
    gen.close()
except RuntimeError as err:
    print(f"error type: {type(err)}")
    print(f"message:    {err}")
\end{verbatim}

Possible output:

\begin{verbatim}
ready
error type: <class 'RuntimeError'>
message:    generator ignored GeneratorExit
\end{verbatim}

A generator may catch \texttt{GeneratorExit} for cleanup, but it must not yield another ordinary value as if execution were continuing.

A safer explicit version is:

\begin{verbatim}
def safe_receiver():
    try:
        yield "ready"
    except GeneratorExit:
        print("closing now")
        raise

gen = safe_receiver()

print(next(gen))
gen.close()
\end{verbatim}

Possible output:

\begin{verbatim}
ready
closing now
\end{verbatim}

The \texttt{raise} statement re-raises the \texttt{GeneratorExit} and allows shutdown to complete.

The control methods can be summarized as:

\begin{center}
\begin{tabular}{|p{0.22\linewidth}|p{0.56\linewidth}|}
\hline
\textbf{Method} & \textbf{Meaning} \\
\hline
\texttt{next(gen)} & resume with \texttt{None}; receive next yielded value \\
\hline
\texttt{gen.send(x)} & resume and make the suspended \texttt{yield} expression evaluate to \texttt{x} \\
\hline
\texttt{gen.throw(exc)} & raise an exception inside the generator at the suspension point \\
\hline
\texttt{gen.close()} & request shutdown by injecting \texttt{GeneratorExit} \\
\hline
\end{tabular}
\end{center}

These are explicit control operations over a suspended frame.

\subsubsection{Cooperative Multitasking: Voluntary Suspension Instead of Preemptive Scheduling}

Generator-based coroutines are cooperative. They run until they reach a suspension point. They do not stop at arbitrary machine instructions.

The suspension point is explicit:

\begin{verbatim}
yield
\end{verbatim}

A simple cooperative task can yield status messages.

\begin{verbatim}
def task(name, count):
    for i in range(1, count + 1):
        print(f"{name}: working step {i}")
        yield f"{name} paused at {i}"

    return f"{name} done"

a = task("A", 3)
b = task("B", 2)

print(next(a))
print(next(b))
print(next(a))
print(next(b))
print(next(a))
\end{verbatim}

Possible output:

\begin{verbatim}
A: working step 1
A paused at 1
B: working step 1
B paused at 1
A: working step 2
A paused at 2
B: working step 2
B paused at 2
A: working step 3
A paused at 3
\end{verbatim}

The caller controls which task resumes next. This is a very small scheduler.

A more systematic scheduler can keep a list of active generators.

\begin{verbatim}
def task(name, count):
    for i in range(1, count + 1):
        yield f"{name}: step {i}"

    return f"{name}: done"

tasks = [
    task("A", 3),
    task("B", 2),
    task("C", 4),
]

while tasks:
    current = tasks.pop(0)

    try:
        message = next(current)
        print(message)
        tasks.append(current)
    except StopIteration as err:
        print(err.value)
\end{verbatim}

Possible output:

\begin{verbatim}
A: step 1
B: step 1
C: step 1
A: step 2
B: step 2
C: step 2
A: step 3
B: done
C: step 3
A: done
C: step 4
C: done
\end{verbatim}

This scheduler performs round-robin execution:

\begin{verbatim}
take first generator
advance it until yield or completion
if it yielded:
    put it at the end of the task list
if it completed:
    remove it permanently
\end{verbatim}

The list of tasks is ordinary.

\begin{verbatim}
def task(name, count):
    for i in range(count):
        yield f"{name}: {i}"

tasks = [
    task("A", 1),
    task("B", 1),
]

print(f"type(tasks): {type(tasks)}")
print()

selected_names = [
    "append",
    "pop",
    "__len__",
    "__getitem__",
    "__iter__",
]

for name in selected_names:
    print(f"{name:<14} {name in dir(tasks)}")

print()
print(f"type(tasks[0]): {type(tasks[0])}")
\end{verbatim}

Possible output:

\begin{verbatim}
type(tasks): <class 'list'>

append         True
pop            True
__len__        True
__getitem__    True
__iter__       True

type(tasks[0]): <class 'generator'>
\end{verbatim}

The scheduler is not using threads. It is not using multiple CPU cores. It is simply choosing which suspended generator frame to resume next.

The important difference from preemptive scheduling is control.

In preemptive scheduling, an external scheduler can interrupt execution without the function explicitly asking to stop.

In generator-based cooperative scheduling, the function must voluntarily yield.

This task cooperates:

\begin{verbatim}
def polite_task():
    print("step 1")
    yield

    print("step 2")
    yield

    print("step 3")
    yield

gen = polite_task()

next(gen)
print("caller regained control")

next(gen)
print("caller regained control")

next(gen)
print("caller regained control")
\end{verbatim}

Possible output:

\begin{verbatim}
step 1
caller regained control
step 2
caller regained control
step 3
caller regained control
\end{verbatim}

This task does not cooperate until the loop finishes.

\begin{verbatim}
def greedy_task():
    for i in range(3):
        print(f"working {i}")

    yield "finally paused"

gen = greedy_task()

print(next(gen))
\end{verbatim}

Possible output:

\begin{verbatim}
working 0
working 1
working 2
finally paused
\end{verbatim}

The caller gets control back only when the generator reaches \texttt{yield}. There is no automatic interruption during the loop.

A cooperative scheduler therefore depends on tasks yielding often enough.

\begin{verbatim}
def fast_task():
    for i in range(3):
        yield f"fast {i}"

def slow_to_yield_task():
    total = 0

    for i in range(1, 5):
        total = total + i
        print(f"slow internal work: {i}")

    yield f"slow total {total}"

tasks = [
    fast_task(),
    slow_to_yield_task(),
]

while tasks:
    current = tasks.pop(0)

    try:
        print(next(current))
        tasks.append(current)
    except StopIteration:
        pass
\end{verbatim}

Possible output:

\begin{verbatim}
fast 0
slow internal work: 1
slow internal work: 2
slow internal work: 3
slow internal work: 4
slow total 10
fast 1
fast 2
\end{verbatim}

The slow task monopolizes execution until it reaches its first \texttt{yield}. Cooperative multitasking requires explicit suspension points.

The architectural model is:

\begin{verbatim}
generator coroutine:
    owns a suspended frame
    runs when caller resumes it
    yields to give control back
    may receive data through send()
    may receive exceptions through throw()
    may be shut down through close()

scheduler:
    chooses which generator to resume
    has control only between yielded suspension points
\end{verbatim}

This is the bridge toward native coroutines. Native \texttt{async def} coroutines use different objects and an event loop, but the conceptual boundary is similar: execution progresses until a structured suspension point, then control returns to an external scheduler.

The final rule is:

\begin{quote}
Generator-based coroutines implement cooperative control flow. They exchange values and control with the caller only at explicit suspension points. No hidden preemptive scheduler interrupts them in the middle of ordinary Python statements.
\end{quote}